% Part: lambda-calculus
% Chapter: introduction
% Section: free-variables

\documentclass[../../../include/open-logic-section]{subfiles}

\begin{document}

\olfileid[pt]{lam}{syn}{fv}
\olsection{Variáveis livres}

O cálculo lambda trata de funções, e a abstração lambda é a maneira como as
funções surgem. Intuitivamente, $\lambd[x][M]$ é a função cujos valores são
dados por~$M$ quando o argumento da função é atribuído a~$x$. Mas nem toda
ocorrência de~$x$ em~$M$ é relevante: se $M$ contém outra abstração
$\lambd[x][N]$, então as ocorrências de $x$ em~$N$ são relevantes para
$\lambd[x][N]$, mas não para $\lambd[x][M]$. Assim, uma abstração lambda
$\lambd[x]$ dentro de $\lambd[x][M]$ \emph{liga} aquelas ocorrências de $x$
em~$M$ que ainda não estão ligadas por outra abstração lambda---as
ocorrências \emph{livres} de~$x$ em~$M$.

\begin{defn}[Escopo]
Se $\lambd[x][M]$ ocorre dentro de um termo~$N$, então a ocorrência
correspondente de~$N$ é o \emph{escopo} de~$\lambd[x]$.
\end{defn}


\begin{defn}[Ocorrência livre e ligada]
  Uma ocorrência da variável $x$ em um termo~$M$ é \emph{livre} se não está
  no escopo de um $\lambd[x]$, e \emph{ligada} caso contrário. Uma
  ocorrência de uma variável~$x$ em $\lambd[x][M]$ é ligada pelo
  $\lambd[x]$ inicial sse a ocorrência de $x$ em~$M$ é livre.
\end{defn}

\begin{ex}
  Em $\lambd[x][x y]$, tanto $x$ quanto $y$ estão no escopo de $\lambd[x]$,
  portanto $x$ é ligada por $\lambd[x]$. Como $y$ não está no escopo de
  nenhum $\lambd[y]$, ela é livre. Em $\lambd[x][x x]$, ambas as ocorrências
  de $x$ são ligadas por~$\lambd[x]$, pois ambas são livres em $xx$. Em
  $((\lambd[x][xx])x)$, a última ocorrência de~$x$ é livre, pois não está no
  escopo de um $\lambd[x]$. Em $\lambd[x][(\lambd[x][x])x]$, o escopo do
  primeiro $\lambd[x]$ é $(\lambd[x][x])x$, e o escopo do segundo
  $\lambd[x]$ é a penúltima ocorrência de~$x$. Em $(\lambd[x][x])x$, a
  última ocorrência de $x$ é livre, e a penúltima é ligada. Assim, a
  penúltima ocorrência de $x$ em $\lambd[x][(\lambd[x][x])x]$ é ligada pelo
  segundo $\lambd[x]$, e a última, pelo primeiro $\lambd[x]$.
\end{ex}

Para um termo $P$, podemos examinar todas as ocorrências de variáveis nele e
obter um conjunto de variáveis livres. Esse conjunto é denotado por
$\FV{P}$, com a seguinte definição natural:

\begin{defn}[Variáveis livres de um termo] \ollabel{def:fv}
  O conjunto das \emph{variáveis livres} de um termo é definido
  indutivamente por:
  \begin{enumerate}
    \item $\FV{x} = \{x\}$ \ollabel{def:fv1} 
    \item $\FV{\lambd[x][N]} = \FV{N} \setminus \{x\}$    \ollabel{def:fv2}
    \item $\FV{PQ} = \FV{P} \cup \FV{Q}$ \ollabel{def:fv3}
  \end{enumerate}
\end{defn}

\begin{prob}
  \begin{enumerate}
  \item Identifique os escopos de $\lambd[g]$ e dos dois $\lambd[x]$ neste
    termo: $\lambd[g][(\lambd[x][g (x x)]) \lambd[x][g (x x)]]$.
  \item Em $\lambd[g][(\lambd[x][g (x x)]) \lambd[x][g (x x)]]$, todas as
    ocorrências de variáveis estão ligadas? Por quais abstrações elas estão
    ligadas, respectivamente?
    \item Dê $\FV{\lambd[x][(\lambd[y][(\lambd[z][x y]) z]) y]}$.
  \end{enumerate}
\end{prob}

\begin{explain}
Uma variável livre é como uma referência ao mundo externo (o
\emph{ambiente}), e um termo que contém variáveis livres pode ser visto como
um termo parcialmente especificado, pois seu comportamento depende de como
configuramos o ambiente. Por exemplo, no termo $\lambd[x][f x]$, que aceita
um argumento~$x$ e devolve $f$ desse argumento, a variável~$f$ é livre. O
valor desse termo depende do ambiente em que ele se encontra, em particular
do valor de $f$ nesse ambiente.

Se aplicarmos abstração a esse termo, obtemos $\lambd[f][\lambd[x][f
    x]]$. Esse termo já não depende da variável de ambiente $f$, pois agora
designa uma função que aceita dois argumentos e devolve o resultado de
aplicar o primeiro ao segundo. Alterar $f$ no ambiente não terá efeito sobre
o comportamento desse termo, pois o termo usará apenas o que lhe for
passado como argumento, e não o valor de $f$ no ambiente.
\end{explain}

\begin{defn}[Termo fechado, combinador]
  Um termo sem variáveis livres é chamado de \emph{termo fechado}, ou
  \emph{combinador}.
\end{defn}

\begin{lem}\ollabel{lem:fv}
  \begin{enumerate}
    \item \ollabel{lem:fv-abs} Se $y \neq x$, então $y \in
      \FV{\lambd[x][N]}$ sse $y \in \FV{N}$.
    \item \ollabel{lem:fv-app} $y \in \FV{PQ}$ sse $y \in \FV{P}$ ou $y
      \in \FV{Q}$.
    \end{enumerate}
\end{lem}

\begin{proof}
  Exercício.
\end{proof}

\begin{prob}
  Prove \olref[lam][syn][fv]{lem:fv}.
\end{prob}

\end{document}

